\documentclass[../../main.tex]{subfiles}% 注意这里的文件路径不能用 ./main.tex ,否则用latexmk编译子文件会报错
\graphicspath{{\subfix{./image/}}{\subfix{../../image/}}} % 指定图片引用路径,后续可以直接使用图片文件名
% 如果图片存放在上级目录,则使用 ../../image/ 作为备用路径

% 例如:
% \begin{figure}[H]
% \centering
% \includegraphics[width=0.8\textwidth]{图.png}
% \caption{}
% \label{figure:图}
% \end{figure}
% 注意:上述\label{}一定要放在\caption{}之后,否则引用图片序号会只会显示??.

\begin{document}

\section{习题}

\begin{example}
设有定义在$E\subseteq \mathbb{R}$上的可测函数列满足
\begin{align*}
f_n(x)\leqslant f_{n+1}(x)\,(n=1,2,\cdots),\ \mathrm{a.e.}\ x\in E.
\end{align*}
若$f_n(x)$在$E$上依测度收敛于$f(x)$，证明：$f_n(x)$在$E$上几乎处处收敛于$f(x)$。
\end{example}
\begin{proof}
由单调有界定理知存在零测集$E_0$，使得对$\forall x\in E\setminus E_0$，有$\lim\limits_{n\to \infty}f_n\left(x\right)\triangleq g\left(x\right)\in \mathbb{R}\cup \left\{+\infty\right\}$。

又由\hyperref[theorem:Riesz定理]{Riesz定理}知，存在子列$\left\{f_{n_k}\right\}$，使得$f_{n_k}\to f,\,\mathrm{a.e.}\,x\in E$。从而存在零测集$E_1$，使得对$\forall x\in E\setminus E_1$，有$\lim\limits_{k\to \infty}f_{n_k}\left(x\right)=f\left(x\right)$。

于是对$\forall x\in E\setminus \left(E_0\cup E_1\right)$，有
\begin{align*}
\lim\limits_{n\to \infty}f_n\left(x\right)=\lim\limits_{k\to \infty}f_{n_k}\left(x\right)=f\left(x\right).
\end{align*}
显然$m\left(E_0\cup E_1\right)=0$，故$f_n\left(x\right)\to f,\,\mathrm{a.e.}\,x\in E$。
\end{proof}







\end{document}